\section{Auxiliary primes and residual concentration}\label{sec:vertical}

The characterwise comparisons of the later sections will invert $p$.
They cannot detect a power of $p$ in a determinant.  The purpose of
this section is to choose finitely many tame variables so that the
relevant determinant already has nonzero reduction modulo $p$.
We do this by making first-order deformations separate every
surviving residual Selmer class.  We give both the arithmetic construction
and the residual-field argument, including reducible representations.

\subsection{The two coefficient families}

Use the fields $K,F,L$ and discriminants $D,D'$ of
Section~\ref{sec:classical}.  In particular $p$ splits completely in
$L$, the residual division field is disjoint from $L$, and both
$E/K$ and $E^{D'}/K$ have rank one and finite Tate--Shafarevich group.
There will be two separate tuples of auxiliary primes.  Over $\Q$
the true variable $t$ is the torsion-free cyclotomic variable, and
this comparison is used only for a non-CM curve.  Over $K$ it is the
anticyclotomic variable, with
\[
 \Gamma\simeq\Z_p
\]
the maximal free pro-$p$ quotient of the ring-class tower at $p$.
The order-class exact sequence shows that inertia at either
$p$-place has open image in $\Gamma$.

Fix supports stable under conjugation.  Over $\Q$ include $p$ and
all bad primes.  For $E/K$ include only places over split rational
primes, including both $p$-places and all bad places.  For the
companion also include the primes dividing $D'$.  Any nonsplit place
in this extra support has no inertia invariant, even on residual
composition factors: $E$ has good reduction there, whereas the
quadratic twist is ramified and $p$ is odd.  In each problem over
$K$, include split good prime ideals whose classes generate the
$p$-part of the ideal class group.  Degree-one representatives exist
by Chebotarev.  Additional fixed split good primes may be included;
we may also prescribe exponents at fixed split test primes without
adding them to the cohomological support.  Outside the support the
condition is unramified.

A tame variable $u_j$ comes from a sequence of distinct good primes
$r_{j,i}$, leaving every finite set and tending $p$-adically to $1$.
Over $K$ these primes split completely in $L$.  Its stage character
has cyclic $p$-power order tending to infinity.  Over $\Q$ it is a
conductor-$r_{j,i}$ cyclotomic character.  Over $K$ it is given by a
surjection
\[
 \lambda_{j,i}:\Pic(\mathcal O_{r_{j,i}})\longrightarrow\Z/p^{n_{j,i}}\Z
\]
which is surjective on inertia at each of the two places over
$r_{j,i}$ and satisfies $\lambda_{j,i}^c=-\lambda_{j,i}$.  A chosen
cyclic generator is represented by $1+u_j$.  Rational characters of
$p$-power order are automatically even.

\begin{lemma}[Auxiliary tuples]\label{vertical:auxiliary-tuples}
For any finite list of fixed split support and test primes, the tame
tuple may be chosen to have the following properties.
\begin{enumerate}
\item Its limiting elliptic Frobenius actions have determinant one
and $V(\gamma_j)-1$ invertible.  Over $K$ the same holds for the
companion.  The weight character $\mathcal C$ used in the theta
construction has nontorsion value on each $\gamma_j$, and its
$\Gamma$ value is nonzero.
\item The vectors of Frobenius exponents in the artificial variables
at the prescribed primes are linearly independent over $\Q_p$,
one vector for each rational prime and with one place chosen in
each split pair.  The list includes the coefficient prime $p$.
\item There are at least two moving primes.  Each moving inertia
group maps onto its own cyclic quotient and trivially to the
others.  All fixed conductors and any prescribed finite set of
primes are avoided.
\end{enumerate}
Here $\mathcal C$ is the algebraic anticyclotomic character with
infinity type one and split conductor used in
\cite[Sections 2 and 4]{C}.  Finitely many such fixed characters may
be treated simultaneously.  In the rational case the conditions on
$\mathcal C$ and $\Gamma$ are omitted.
\end{lemma}

\begin{proof}
First prescribe power-residue data.  In the rational case the
radicands are the fixed primes themselves.  In the imaginary case
choose ideal-class generators away from the prescribed places,
generators of their principal power relations, and generators for
expressions of the prescribed prime ideals in those classes.  Include
the conjugates.  For a split prime $\mathfrak q$, take a class-number
power $\mathfrak q^h=(a_{\mathfrak q})$.  The corresponding
anticyclotomic residue exponent is, up to its common generator
normalization, the Kummer exponent of
$a_{\mathfrak q}/\overline{a_{\mathfrak q}}$.  These elements are
independent by their valuations at the chosen split pairs.

Their Kummer image over the cyclotomic tower has full $\Z_p$-rank.
One way to see the required restriction injectivity modulo $p^a$
is to use a central cyclotomic element whose scalar differs from
$1$ by a unit; such an element exists since $p$ is odd and splits
in the fixed fields.  A further fixed finite extension loses at
most a bounded power of $p$.  A $\Z_p$ relation annihilating the
Kummer image would consequently give power relations modulo
arbitrarily high $p$-powers, contradicting the independent
valuations.  Thus sufficiently deep prescribed exponents can be
chosen independently.

To extend an inertia logarithm across the class group, write the
fixed class-group relations in invariant-factor form.  Require the
power-residue exponents to be divisible by the fixed $p$-parts of
these invariant factors, and divide the relations to assign the
values on the ideal generators.  Principal relations then define
the whole character.  Anti-invariance is imposed using the
conjugate generators.  The fixed divisions change neither the full
rational span of the limiting exponents nor their independence.
Finite Chebotarev conditions on roots of unity and on this radical
list realize these power-residue prescriptions.  Taking enough
moving primes gives the independent row vectors in the assertion.

We next check compatibility with the elliptic and fixed character
conditions.  Adjoin all $p$-power roots of elements of the base to
the cyclotomic tower, and consider the joint compact $p$-adic
analytic image of the Tate representation and the fixed characters.
The common Lie quotient with the radical translation group is
abelian.  In the non-CM case it kills the $\mathfrak{sl}_2$ factor
of the Tate image, by Serre's open image theorem \cite{Serre}.
In the CM case, after a fixed extension, the connected Tate image
is a torus.  Conjugation on the remaining abelian analytic quotient
is trivial after a finite extension, whereas conjugation on radical
translations is multiplication by the cyclotomic character.
Choosing that character different from $1$ proves that the common
Lie quotient is zero.  Hence fixing all radicals removes no Lie
dimension from the analytic image.

Make the prescribed radical data sufficiently deep that an identity
neighborhood remains available.  Its determinant-one Tate projection
contains an open special linear subgroup, or an open norm-one CM
torus.  The equation $\det(V(\gamma)-1)=0$ is a proper analytic
condition there.  The weight character is nonconstant by its
nonparallel labelled $p$-adic exponents, and the true
anticyclotomic character is nonconstant by its inertia description.
After shrinking the neighborhood, the torsion condition on the
weight is the zero-logarithm locus.  Avoiding these finitely many
proper analytic loci gives the required Frobenius limits.  The
companion has the same Tate action on $G_L$.

Only finite radical and coefficient tables are used at any stage.
They can therefore be combined into finite Chebotarev prescriptions
with accuracy tending to infinity.  If a bound on
$v_p(r_{j,i}-1)-n_{j,i}$ is required, perturb by a sufficiently high
power of a fixed near-identity element with nontorsion cyclotomic
projection.  Its radical geometric sums vanish to the required
precision, while its cyclotomic displacement has valuation
$n_{j,i}+O(1)$.
No boundedness of the unramified $p$-local degrees of the full
conductor fields is used.
\end{proof}

Put $R=\Z_p[[t,\mathbf u]]$ and let $\chi$ be the product of the
true and artificial character scalars.  Over $\Q$ let $G$ be the
ordinary global complex for $T\chi$, with the moving places allowed.
Over $K$ use the full global complex, also allowing the moving
places, and set
\[
 C_w=\operatorname{fib}(G\longrightarrow L_w^{\rm loc}),
 \qquad w\mid p,
\]
where the condition at the other $p$-place is full.  The residual
restriction kills global $H^0$, and the complementary restriction
at the other $p$-place kills dual $H^0$.  Euler characteristic is
zero.  Thus $C_w$ has a square free model in degrees $1,2$; let
$L_w$ denote its differential determinant, up to a unit.  Use
$L'_w$ for the companion.  Inverting all character conventions
gives the same assertions.

\begin{proposition}[Residual concentration]\label{vertical:tests}
After adjoining finitely many more tame variables, the separate
rational and imaginary tuples satisfy
\begin{equation}\label{eq:source-1}\tag{1}
 \begin{array}{ll}
 \mathrm{(a)}&G\otimes_RR_{(p)}\text{ is a free line in degree one};\\[2pt]
 \mathrm{(b)}&L_w,L'_w\not\equiv0\pmod{(p,t)}
                  \quad\text{for both }w\mid p.
 \end{array}
\end{equation}
\end{proposition}

The proof occupies the rest of the section.  We first describe the
residual cohomology and its parity, then produce a deformation that
kills each residual strict space.

\subsection{Residual spaces and the scalar primitive}

Initially specialize to
\[
 k_r=\Frac(\F_p[[\mathbf u]]),\qquad t=0.
\]
The residual elliptic representation is either absolutely simple or
has character constituents.  Indeed its image contains an element
acting with eigenvalues $1,-1$, even after restriction to $G_K$,
by the disjointness imposed in Section~\ref{sec:classical}.  An
irreducible plane cannot have a quadratic endomorphism field
commuting with this element.  In the reducible case the two
characters are distinct, also over $K$.  Test each simple
constituent, extending the finite coefficient field when necessary,
and include the $\chi$ twist.  Its paired Tate dual has the inverse
twist.

At every finite place in these tests $H^0=H^2=0$.  At fixed split
places the nonconstant Frobenius scalar gives $H^0=0$; at nonsplit
exceptional places inertia does so; at a moving place its own
inertia variable does so.  Apply the same argument to the dual to
obtain $H^2=0$.  Away from $p$, Euler characteristic zero then
makes the full local complex acyclic.  Global invariants vanish
as well.  The strict Selmer degree-one spaces consequently inject
into unrestricted cohomology by Lemma~\ref{models:evaluations}.

Over $\Q$, write $s$ for the dimension of the dual all-strict
space.  Its paired unrestricted primal space has dimension $s+a$,
where $a$ is the negative dimension of the primal constituent at
infinity.  Killing the former space will kill primal $H^2$.
Over $K$, pair the $w$-strict space with the dual $\bar w$-strict
space.  Exact orthogonality and Euler characteristic make their
dimensions equal, say $s$.

\begin{lemma}[Scalar localization]\label{vertical:scalar-localization}
For the constant roots coefficient, global degree-two localization
is injective on these residual tests.  In particular a cup of two
paired classes which is locally zero has a global cochain primitive.
\end{lemma}

\begin{proof}
At residual arithmetic precision, the dual of the localization
kernel is the space of everywhere locally trivial constant
$\F_p$-characters, by Poitou--Tate.  Such a character is unramified
everywhere and split at the support; it is therefore a class-group
character killed by the chosen generating prime classes.  Over
$\Q$ the same group is zero.  The scalar representation and its
cohomological comparison do not depend on the artificial variables.
Tensoring its residual diagram with $\F_p[[\mathbf u]]$ and then
its fraction field preserves this injection.  Equivalently retain
the injectivity cone in Lemma~\ref{models:finite-models}.  Thus a
locally zero degree-two scalar class is zero, giving the primitive.
\end{proof}

\begin{lemma}[Conjugate parity]\label{vertical:parity}
For an absolutely simple residual elliptic plane over $K$, each
one-sided strict space above has even dimension.  The same
alternation statement holds on characteristic-zero localizations
with unramified conditions at fixed non-$p$ places and constant
inertia action.
\end{lemma}

\begin{proof}
Keep $t=0$ and lift the residual test to the $(p)$-DVR of
$\Z_p[[\mathbf u]]$.  Non-$p$ local complexes are contractible
there, since their residual tests are acyclic.  The remaining
one-sided problem is perfect in degrees $1,2$, with conjugate
self-duality.  It is generically acyclic.  To check this last
assertion, specialize centrally in characteristic zero.  The
moving-place singular determinants are nonzero because
$V(\gamma_j)-1$ is invertible.  Inflation therefore gives the
fixed-support problem.  The simple-zero choices give one compact
Selmer line from points, and its localization at either $p$-place
is nonzero because the local logarithm has torsion kernel.  At
non-$p$ places the rational local elliptic cohomology vanishes by
the absence of elliptic invariants, local duality, and Euler
characteristic zero.  The strict--finite triangles of
Lemma~\ref{models:strict-finite} now give strict acyclicity.  A
nonzero central determinant is nonzero generically.

We spell out the sign that controls specialization.  On the induced
coefficient, in its two $K$-coordinates, the pairing is
\begin{equation}\label{vertical:induced-pairing}
 ((a,b),(a',b'))\longmapsto e(a,b')+e(b,a'),
\end{equation}
where $e$ is the elliptic alternating pairing.  This is alternating
and Tate-equivariant, including under complex conjugation, which
swaps the coordinates and reverses $e$.  Its compact cup gives a
graded skew pairing of shift $-3$.  Split-coordinate local
conditions are isotropic.  Unramified conditions have the residue
cup nullhomotopy, and contractible local conditions may be replaced
by zero.  At a newly allowed discriminant place the Tate module is
unramified over $K$; quadratic inertia on the induced coefficient
swaps the two coordinates.  Its invariants form the graph of a
character unit.  Restriction to this graph has the same residue
nullhomotopy, since the residue group has cohomological dimension
one and the quadratic index is invertible.  The real Tate terms
are zero.

Graded cup commutativity, the cone cup, and their homotopies are
identities at finite $p$-power precision.  Transfer them by
Lemma~\ref{models:finite-models}.  Antisymmetrization is allowed
because $2$ is a unit.  On a minimal two-term model the resulting
perfect pairing between degrees one and two identifies the
differential with an alternating matrix: if $P$ denotes the
degree-three pairing, its chain identity gives
$P(x,dy)=P(dx,y)=-P(y,dx)$.  The residual and generic nullities of
an alternating matrix have the same parity.  Generic acyclicity
therefore gives even residual nullity.

This also explains the precise extension of the conjugate
Bockstein argument in \cite[Lemma 12.1 and Corollary 12.2]{B}.
No characteristic-zero division is used in
transferring the pairing to the residual field.  For the later
characteristic-zero localizations the same calculation applies
with the constant-inertia unramified comparisons of
Section~\ref{sec:models}; their nullhomotopies are again the
residue ones.
\end{proof}

\subsection{A new tame direction and its lifting obstruction}

Adjoin a new split tame prime sequence with logarithm $\lambda$.
Require $\lambda$ to vanish at all old places, including the old
moving places, at every stage.  This is achieved by imposing
sufficient radical precision on their ideal relations.  Require
also that the old artificial character on its limiting Frobenius
$\gamma$ have a nonzero exponent.  When the new variable is set
to zero, its local complex is then acyclic over $k_r$; unramified
inflation identifies the enlarged specialization with the old
diagram.

These requirements are compatible with prescribed old evaluations.
Start with an element fixing all radicals from the base and having
the Tate and fixed-character conditions of
Lemma~\ref{vertical:auxiliary-tuples}.  Adjust the old scalar
exponents by old moving-prime inertia.  Such inertia is trivial
on the full Tate module, on the fixed splitting fields and
characters, and on the true anticyclotomic character.  To make
an adjustment fix all radicals too, multiply an inertia element
by its conjugate under an element over the base acting by $-1$
on $\mu_{p^\infty}$.  The radical translations cancel whereas
the old abelian character exponents double.  This factor of two
is invertible.  Existence of the conjugating element follows
from the splitting of $p$ in the fixed fields.  Make these
adjustments term by term, retaining the finite evaluation tables.

\begin{lemma}[Necessary first-order lift test]\label{vertical:lift-test}
Let $x$ be an old strict class and $d$ a class in its paired opposite
space.  Choose a primitive $u$ of their scalar cup,
$du=x\cup d$.  For whole spaces of such classes, fix cochain
representatives on bases, choose the primitives on pairs of basis
elements, and extend these choices bilinearly.  Suppose $x$ lifts
in the new direction, with
derivative satisfying $dx_1=-\lambda\cup x$ in the chosen character
convention.  Put
\[
 A=(M_x(\gamma)-1)^{-1}x(\gamma),
\]
where $M_x$ is the coefficient representation of $x$.  Then the
Frobenius evaluation $z_\gamma$ of $u-A\cup d$ vanishes, up to a
common nonzero sign or coordinate factor.  Over $K$ the necessary
equation is its difference with the conjugate evaluation, since
the conjugate tame logarithm has the opposite sign.  The test is
linear in the new direction and bilinear in $x,d$, and is valid
over $k_r$ as well as over the characteristic-zero field tests.
\end{lemma}

\begin{proof}
The primitive exists by Lemma~\ref{vertical:scalar-localization}.
At each $p$-place one of the paired classes is strict, with the
specified local nullhomotopy; away from $p$ the residual local
complexes are acyclic.  At the later horizontal tests the paired
unramified cups instead vanish through the residue complex of
cohomological dimension one.  Thus the cup is locally zero.
Apply the
arithmetic sum of local invariants to the degree-two cochain
\begin{equation}\label{vertical:primitive-cup}
 x_1\cup d-\lambda\cup u,
\end{equation}
with the simultaneous sign convention making it closed.  At old
places $\lambda$ is zero and the local condition removes the
remaining cup.  At the new place $x,d,u$ are unramified.  The
Frobenius-minus-one operator is invertible, and $A$ trivializes
$x$ in the residue complex.  Subtracting this coboundary leaves
the unramified scalar cochain $u-A\cup d$.  Cup with the tame
direction has invariant equal, up to sign, to its Frobenius value.
Here the Frobenius fixes the $p$-power roots of unity, and the
direction is normalized by the power-residue exponent in the
compatible root basis.  The mixed residue--inertia cup is therefore
the usual local reciprocity pairing.  Over $K$ the two places
give the indicated difference.  If several new directions are
combined, the logarithm of one is unramified at the other new
places and gives no additional mixed contribution.

We verify the residual interpretation explicitly.  Mark the old
global cochain inclusions and their values on the Frobenius
sequences, then use the inflation homotopies in adjoining the
new primes.  Their restrictions to new inertia are zero, and
the induced residue representatives have the Frobenius values
just prescribed.  The invertible residue operator removes their
degree-one cohomology, leaving exactly the scalar residue
cochain above.  Mark its root basis and the mixed cup map too.
All the coboundary corrections, restriction homotopies, and the
compact invariant equation are finite matrix identities from
actual arithmetic diagrams.

To use classes over $k_r$, clear their finitely many denominators
in $\F_p[[\mathbf u]]$, including those in the solution for $u$
and in the inverse residue operator.  The resulting equations
hold in the corresponding localized finite free models.  Any
equation which holds only after that localization becomes an
equation over $\F_p[[\mathbf u]]$ after multiplication by one
more common nonzero parameter denominator.  Now use the Artin
rings $A_b=\F_p[[\mathbf u]]/(\mathbf u)^b$, with the first-order
new variables adjoined as needed.  The prime $p$ is zero
throughout this residual calculation.  The cochain inclusions
and homotopies identify its equations with the corresponding
actual stage arithmetic identities over $A_b$; no lift of an
arbitrary residual solution to $\Z/p^n$ is asserted.  Passing
to the limit in $b$ and then inverting the same
nonzero parameter denominators proves the stated equation over
$k_r$.  This does not divide by $p$.  It also shows that further
field extensions preserve the test.  For characteristic-zero
field points one may instead clear fixed scalar denominators
and increase the arithmetic precision, as in
Lemma~\ref{models:finite-models}.  Thus the primitive computation applies
with its residual arithmetic identities established explicitly.
\end{proof}

\subsection{Spanning the obstruction matrices}

Choose bases $x_1,\ldots,x_s$ of a strict space and
$d_1,\ldots,d_{s+a}$ of the paired opposite space in the rational
case, or $d_1,\ldots,d_s$ in the imaginary case.  For a new
direction use the fixed pairwise primitives $u_{ij}$ with
$du_{ij}=x_i\cup d_j$, extended bilinearly as above.  Its
\emph{necessary lift matrix} at Frobenius $\gamma$
has $(i,j)$-entry $z_\gamma(x_i,d_j)$ from
Lemma~\ref{vertical:lift-test}, or the conjugate difference there
over $K$.  A lifted class with row of coefficients $\alpha$ must
satisfy $\alpha Z_\gamma=0$.  Thus a matrix with zero left kernel
excludes every nonzero lift in that direction.

\begin{lemma}[Obstruction span]\label{vertical:obstruction-span}
For the residual tests above, finitely many new tame directions
have necessary lift matrices whose linear span contains every
$s\times(s+a)$ matrix over $k_r$ in the rational problem, every
$s\times s$ matrix for a reducible constituent in the imaginary
problem, and every alternating $s\times s$ matrix for an
absolutely simple plane in the imaginary problem.  Previously
prescribed Tate actions, radical data, old character actions,
and acyclicity of the new local complexes can all be retained.
\end{lemma}

\begin{proof}
Let $J$ be the termwise joint kernel of the full Tate and old
character actions, including $L$ and the fixed splitting data
over $K$.  It is harmless to retain the fixed character data
as well.  On $J$ the evaluations of $x,d$ are additive.
They detect their classes: old inertia at one moving prime
gives a central element in the quotient image whose scalar
difference from $1$ is invertible.  Apply
Lemma~\ref{models:evaluations}.  Independent classes in a simple
representation have evaluations spanning their full tuple,
jointly for nonisomorphic representations.

In the rational case the opposite representations are already
nonisomorphic on old inertia, which sees $\chi$ and $\chi^{-1}$.
In the imaginary reducible case include the conjugate evaluations.
The resulting four representations are distinct: the elliptic
characters are distinct, and opposite infinite scalar twists
cannot be identified.  For an absolutely simple plane take
$d_j=x_j^c$ under conjugate duality.  The two oppositely twisted
planes are nonisomorphic for the same inertia reason.

Multiply a chosen Frobenius representative by a commutator
$[v,w']$ with $v,w'\in J$.  Since the full Tate action is fixed,
both elements fix the cyclotomic tower.  Their radical actions
are translations in an abelian group.  The commutator therefore
fixes all radicals, as well as the Tate and old character data.
The additive evaluations of $x,d$ on the commutator are zero,
so the local correction involving $A$ is unchanged.  The
primitive equation changes its scalar evaluation by
\begin{equation}\label{vertical:commutator-change}
 -e(x(v),d(w'))+e(x(w'),d(v)),
\end{equation}
up to the fixed sign convention.

Joint spanning and bilinearity now give elementary rectangular
matrices in the rational problem and elementary square matrices
in a character problem.  For a plane over $K$, take the
conjugate difference.  With $d_j=x_j^c$, the transposed
comparison has the opposite sign by Weil skew-symmetry; hence
the changes span $E_{ij}-E_{ji}$, the elementary alternating
matrices.  This uses linear span over $k_r$, not an assertion
that every vector in that span is a single Galois evaluation.
Only finitely many commutators are needed to span the finite
matrix spaces.  Differences from the initial Frobenius test
belong to the span of the resulting directions, so their
linear span has the asserted size.

Realize the required Frobenius representatives by Chebotarev at
increasing finite precision, preserving the radical and
evaluation tables of Lemma~\ref{vertical:lift-test}.  Normalize
all root and tangent bases in the same diagrams.  A common
nonzero sign, or a nonzero scalar multiplying an entire tangent
direction, does not change its span.  For conjugate places use
the transported marking, so the two entries are compared with
the same normalization.
\end{proof}

\begin{proof}[Proof of Proposition~\ref{vertical:tests}]
If a tested strict space is zero no new direction is needed.
Otherwise Lemma~\ref{vertical:obstruction-span} supplies a
linear combination of directions whose necessary lift test has
zero left kernel.  In the plane case this uses
Lemma~\ref{vertical:parity}: $s$ is even and an alternating
matrix with $s/2$ invertible skew blocks is nonsingular.

Work in the enlarged residual ring, localized where the old
variables have fraction field $k_r$ and all new variables are
zero.  The new local complexes are acyclic at this point, so
the specialized global diagrams are the old ones.  Cancel their
unit differential blocks.  In the remaining minimal complex,
the derivative of a lifted degree-one vector in a new tangent
direction is the first-order differential.  A vector in its
kernel would pass every necessary cup test.  The chosen test
has zero kernel, so this first-order differential has full
column rank in that direction.  Some maximal minor of the
actual differential has a nonzero initial term.  It is
therefore nonzero over the enlarged generic residual field,
which kills that strict space.

Take the union of the finite direction lists for all simple
constituents, both sides, and both imaginary curves.  Each test
retains its own slice, since every new local prime was chosen
Frobenius-acyclic there.  Long exact sequences for the residual
composition filtration then give the desired vanishing for
the original representation, including nonsplit extensions.

Over $\Q$ the dual strict vanishing gives $H^2=0$ for the
primal full complex.  Its Euler characteristic is $-1$, so it
has one residual degree-one line.  Minimality first at
$(p,t)$ and then at $(p)$ gives a free line over $R_{(p)}$.
Over $K$, one-sided $H^1$ and $H^2$ vanish at the generic
residual test.  The square determinant consequently has
nonzero reduction modulo $(p,t)$, on both sides and for the
companion.  This is~\eqref{eq:source-1}.
\end{proof}

\subsection{Further directions and characteristic-zero tests}

The same construction will be needed when removing horizontal
divisors.  We record what can be retained, so that the later
commutative algebra need not repeat the arithmetic choices.

Here a \emph{horizontal field test} means the derived fiber over
$k(\mathfrak P)=\Frac(R/\mathfrak P)$ for a height-one prime
$\mathfrak P$ of $R=\Z_p[[t,\mathbf u]]$ with $p\notin\mathfrak P$.
Its local testing ring is the DVR $R_{\mathfrak P}$.  This field
need not be a finite extension of $\Q_p$.  The finite-extension
interior points of Lemma~\ref{models:field-points} are points of
the divisor at which finitely many specified denominators remain
nonzero.  Comparisons then specialize from the localization of
$R/\mathfrak P$ at those denominators, not from its entire
fraction field.  In an enlargement by new variables $\mathbf v$, retaining
an old test means localization at $(\mathfrak P,\mathbf v)$ and
specialization $\mathbf v=0$ before its field fiber is taken.

\begin{lemma}[Preservation under enlargement]\label{vertical:enlargement}
Property~\eqref{eq:source-1}(b) persists when further tame primes
are adjoined, provided each has a nonzero Frobenius exponent
in the preceding artificial variables.  On setting the new
variables to zero, unramified inflation identifies the old
diagram, with the full new-place singular factors retained.
Those factors have nonzero reduction on the residual test.
\end{lemma}

\begin{proof}
At the old residual test the new prime has unramified scalar
Frobenius involving an old variable with nonzero exponent.
Its Frobenius-minus-one determinants are nonzero; the local
complex is acyclic over the residual fraction field.  The
localized full-to-unramified triangle multiplies the
determinant by a nonzero residual singular factor.  Hence
specialization of each enlarged determinant remains nonzero
at the old residual test, which proves the assertion.
\end{proof}

\begin{lemma}[Directions at horizontal tests]\label{vertical:horizontal-directions}
Fix finitely many characteristic-zero height-one field tests
of a tuple satisfying~\eqref{eq:source-1}(b), with at least two
old artificial variables.  One may adjoin finitely many tame
directions, locally zero on all old support, so that the new
prime complexes are acyclic at all these fields and the
necessary lift tests of Lemma~\ref{vertical:lift-test} retain
their commutator-spanning property.  The construction is
compatible with further field extension and with the
primitive unramified conditions at fixed non-$p$ places.
\end{lemma}

\begin{proof}
At a height-one test at least one old inertia scalar has
infinite order: two independent variables cannot both
specialize to roots of unity on a single prime divisor.
Adjust the old Frobenius exponents of a new prime as above.
On a line in an infinite-order coordinate, an elliptic
eigenvalue exception excludes only boundedly many integer
choices.  A sufficiently large finite grid in all the old
exponents therefore avoids the exceptions for the finite
list of fields and both representations.  Keep the
nonzero-exponent requirement, all Tate and fixed character
conditions, and the radical conditions making the new
logarithms zero at old places.  This gives simultaneous
Frobenius acyclicity.

Exact orthogonality at the old unramified places uses
Lemma~\ref{models:orthogonality}.  Their cup contributions
and first-order contributions vanish through its residue
nullhomotopy, since the new logarithm is locally zero.
Scalar localization is obtained from the constant model
and Lemma~\ref{vertical:scalar-localization}; it remains
injective after these characteristic-zero specializations.
The marked cup and primitive proof now gives the lift
identity at each field.  On the joint kernel, evaluations
are detected by the central old inertia scalar, and the
commutator formula~\eqref{vertical:commutator-change}
preserves all the imposed data.  For any finite list of
fields take the union of their direction lists.  Each
retains its own spanning slice; it is unnecessary to
prescribe independent matrices at all fields with a single
prime.  Clearing the finitely many denominators and slowing
precision preserves the comparison identities, also after
field extension.
\end{proof}

The vertical part of the comparison is now fixed integrally.
The later sections may test horizontal divisibility after
inversion of $p$ while retaining the nonzero residual
determinants in~\eqref{eq:source-1}.
